% ============================================================
%  Lecture 5 — C++ recap & software architecture for optimisation
% ============================================================

\section{C++ Recap}
\label{sec:cpp-recap}

From this point on we \emph{implement} the algorithms we study.
The language is C++ (compiled with \texttt{g++} or
\texttt{clang++}).
This section recaps the OOP features we rely on; students
comfortable with all of them can skip ahead to
\cref{sec:solver-architecture}.

% ------------------------------------------------------------
\subsection{Encapsulation and information hiding}
\label{ssec:encapsulation}
% ------------------------------------------------------------

A \textbf{class} bundles \emph{data members} (the state) and
\emph{methods} (the operations on that state) into a single unit.
Access control enforces \emph{information hiding}:

\begin{itemize}[nosep]
  \item \texttt{private} (default for \texttt{class}):
    accessible only within the class itself.
  \item \texttt{protected}: accessible within the class and its
    subclasses.
  \item \texttt{public}: accessible from anywhere.
\end{itemize}

\begin{lstlisting}[language=C++, caption={Basic class layout.}]
class Graph {
public:                     // interface
  Graph(string file_name);  // constructor
  unsigned Nodes() const { return nodes; }
  bool Arc(unsigned i, unsigned j) const;
protected:                  // internal state
  unsigned nodes;
  vector<vector<bool>> adjacency_matrix;
};
\end{lstlisting}

\noindent
The \texttt{const} qualifier on \texttt{Nodes()} promises that
calling it does not modify the object---a crucial contract for
readers of the code and for the compiler's optimiser.

% ------------------------------------------------------------
\subsection{Constructors}
\label{ssec:constructors}
% ------------------------------------------------------------

A \textbf{constructor} initialises an object when it is created.
It has the same name as the class and no return type.
For our optimisation code, the constructor typically reads an
instance from a file:

\begin{lstlisting}[language=C++, caption={Constructor that reads
  from a file.}]
Graph::Graph(string file_name) {
  ifstream is(file_name);
  if (!is) {
    cerr << "Could not open " << file_name << endl;
    exit(1);
  }
  // ... parse the file and fill data members ...
}
\end{lstlisting}

% ------------------------------------------------------------
\subsection{File I/O: \texttt{ifstream} and \texttt{ofstream}}
\label{ssec:file-io}
% ------------------------------------------------------------

The standard library provides \texttt{ifstream} (input) and
\texttt{ofstream} (output) in \texttt{<fstream>}.
They behave like \texttt{cin}/\texttt{cout} but read from / write
to files:

\begin{lstlisting}[language=C++, caption={Reading and writing files.}]
#include <fstream>

ifstream is("instance.txt");   // open for reading
unsigned n;
is >> n;                       // read an integer

ofstream os("solution.txt");   // open for writing
os << "makespan = " << cmax << endl;
\end{lstlisting}

% ------------------------------------------------------------
\subsection{The \texttt{vector} container}
\label{ssec:vector}
% ------------------------------------------------------------

\texttt{std::vector<T>} is a dynamically-sized array.  We use it
everywhere instead of raw C arrays:

\begin{lstlisting}[language=C++, caption={Common \texttt{vector}
  operations.}]
vector<int> v(10, 0);        // 10 elements, all zero
v.push_back(42);             // append
v.size();                    // current length
v[3] = 7;                    // random access (no bounds check)

// 2D: vector of vectors
vector<vector<bool>> matrix(n, vector<bool>(n, false));
\end{lstlisting}

\noindent
For adjacency matrices we use
\texttt{vector<vector<bool>{}>} (compact, random access);
for adjacency lists, \texttt{vector<vector<unsigned>{}>}
(compact, iteration-friendly).


% ============================================================
\section{Software Architecture for Optimisation}
\label{sec:solver-architecture}
% ============================================================

Every solver we build in this course follows the same three-module
architecture:%
\aside{Separating I/O from the solver
lets us swap algorithms without
touching file-parsing code.}

\begin{figure}[H]
\centering
\begin{tikzpicture}[
    box/.style={draw=blue!60!black, fill=blue!6,
                rounded corners=3pt, minimum width=28mm,
                minimum height=12mm, font=\small\bfseries,
                align=center},
    arr/.style={-Stealth, thick, blue!60!black},
    x=42mm, y=16mm
  ]
  \node[box] (inp) at (0, 0) {Input\\(\texttt{Data})};
  \node[box] (sol) at (1, 0) {Solver};
  \node[box] (out) at (2, 0) {Output\\(\texttt{Solution})};

  \draw[arr] (inp) -- (sol) node[midway, above,
    font=\footnotesize] {reads};
  \draw[arr] (sol) -- (out) node[midway, above,
    font=\footnotesize] {writes};

  % file labels
  \node[font=\footnotesize\itshape, text=gray, below=4mm of inp]
    {instance file};
  \node[font=\footnotesize\itshape, text=gray, below=4mm of out]
    {solution file};
\end{tikzpicture}
\caption{Three-module architecture.  The \textbf{Input} class
         reads the instance file and stores the problem data.
         The \textbf{Solver} reads from Input and writes into
         a \textbf{Solution} object.  The Solution class handles
         output formatting.}
\label{fig:three-module}
\end{figure}

\begin{enumerate}
  \item \textbf{Input / Data class.}
    Reads and stores the instance.  Provides read-only accessors
    (\texttt{const} methods) so that the solver cannot
    accidentally modify the input.

  \item \textbf{Output / Solution class.}
    Stores the current solution and provides operators for
    modifying and printing it.

  \item \textbf{Solver.}
    Contains the algorithm logic.  It receives a \texttt{const}
    reference to the Input and produces a Solution.
    Swapping solvers (greedy, backtracking, local search) requires
    no changes to the I/O classes.
\end{enumerate}

This separation is not just good software engineering---it is
essential when comparing multiple algorithms on the same
instances, which is exactly what the course project requires.


% ============================================================
\section{Example: Graph Colouring I/O Classes}
\label{sec:gc-data}
% ============================================================

We illustrate the architecture on the graph colouring problem.
The two files \texttt{GC\_Data.hh} (header) and
\texttt{GC\_Data.cc} (implementation) define the Input and Output
classes.

% ------------------------------------------------------------
\subsection{The \texttt{Graph} class (Input)}
\label{ssec:graph-class}
% ------------------------------------------------------------

\begin{lstlisting}[language=C++,
  caption={\texttt{GC\_Data.hh} --- the \texttt{Graph} class
           (input data).},
  label={lst:gc-data-hh-graph}]
class Graph {
  friend ostream& operator<<(ostream& os, const Graph& g);
public:
  Graph(string file_name);
  unsigned Nodes()   const { return nodes; }
  unsigned Colors()  const { return colors; }
  bool Arc(unsigned i, unsigned j) const
    { return adjacency_matrix[i][j]; }
  unsigned Degree(unsigned i) const
    { return adjacency_list[i].size(); }
  unsigned ArcListMember(unsigned i, unsigned j) const
    { return adjacency_list[i][j]; }
protected:
  unsigned nodes, colors;
  vector<vector<bool>>     adjacency_matrix;
  vector<vector<unsigned>> adjacency_list;
};
\end{lstlisting}

\paragraph{Design notes.}
\begin{itemize}[nosep]
  \item The graph is stored in \emph{two} representations:
    an adjacency matrix (fast $\Oh{1}$ edge queries) and
    adjacency lists (fast neighbour iteration).
    Both are built once by the constructor.
  \item All accessors are \texttt{const}---the solver can read but
    never modify the graph.
  \item The \texttt{friend} declaration lets the \texttt{<\!<}
    operator access private members for printing.
  \item Data members are \texttt{protected} (not \texttt{private})
    so that subclasses---e.g.\ a \texttt{WeightedGraph}---can
    extend them.
\end{itemize}

\paragraph{The constructor} reads a MiniZinc-style data file:

\begin{lstlisting}[language=C++,
  caption={Constructor: reading the instance file.},
  label={lst:gc-constructor}]
Graph::Graph(string file_name) {
  char ch; string buffer; unsigned i, j, val;
  ifstream is(file_name);
  if (!is) { cerr << "Could not open " << file_name << endl; exit(1); }

  is >> buffer >> ch >> nodes  >> ch;  // "nodes = <n>;"
  is >> buffer >> ch >> colors >> ch;  // "colors = <k>;"

  adjacency_matrix.resize(nodes, vector<bool>(nodes, false));
  adjacency_list.resize(nodes);

  is >> buffer >> ch >> ch >> ch;      // "adjacency_matrix = [|"
  for (i = 0; i < nodes; i++)
    for (j = 0; j < nodes; j++) {
      is >> val >> ch;
      if (val == 1 && i < j) {        // upper triangle only
        adjacency_matrix[i][j] = true;
        adjacency_matrix[j][i] = true;
        adjacency_list[i].push_back(j);
        adjacency_list[j].push_back(i);
      }
    }
}
\end{lstlisting}

\noindent
The file format matches MiniZinc's data syntax (which we will
use in the next lectures), so the same instance file serves both
the C++ solver and the MiniZinc model.

% ------------------------------------------------------------
\subsection{The \texttt{Coloring} class (Output)}
\label{ssec:coloring-class}
% ------------------------------------------------------------

\begin{lstlisting}[language=C++,
  caption={\texttt{GC\_Data.hh} --- the \texttt{Coloring} class
           (solution).},
  label={lst:gc-data-hh-coloring}]
class Coloring {
  friend ostream& operator<<(ostream& os, const Coloring& col);
public:
  Coloring(const Graph& my_in)
    : in(my_in), coloring(in.Nodes(), -1) {}
  void Reset();
  int  operator[](unsigned i) const { return coloring[i]; }
  int& operator[](unsigned i)       { return coloring[i]; }
protected:
  const Graph& in;
  vector<int>  coloring;
};
\end{lstlisting}

\paragraph{Design notes.}
\begin{itemize}[nosep]
  \item The constructor takes a \texttt{const Graph\&} and
    initialises every node's colour to $-1$ (uncoloured).
  \item The \texttt{operator[]} overloads provide array-like
    syntax: \texttt{col[v] = 3} assigns colour~3 to node~$v$;
    \texttt{col[v]} reads it.
    The \texttt{const} version returns by value (read-only);
    the non-\texttt{const} version returns by reference
    (read-write).
  \item \texttt{Reset()} sets all colours back to $-1$, allowing
    the solver to restart without reallocating memory.
  \item The \texttt{in} reference lets the output class access
    graph data (e.g.\ number of nodes) without duplicating it.
\end{itemize}

\begin{observation}
This pattern---\texttt{Data} class with \texttt{const} accessors,
\texttt{Solution} class with a reference to the data---recurs in
every problem we implement.  For TSP: \texttt{TSP\_Data} stores
the distance matrix, \texttt{TSP\_Solution} stores the tour.
For bin packing: \texttt{BP\_Data} stores item sizes and bin
capacity, \texttt{BP\_Solution} stores the assignment of items to
bins.  The solver is always a separate function or class that
reads from \texttt{Data} and writes to \texttt{Solution}.
\end{observation}


% ============================================================
\section{Implementing Greedy: DSatur for Graph Colouring}
\label{sec:gc-greedy-impl}
% ============================================================

With the architecture in place, let us implement the DSatur
algorithm from \cref{ssec:dsatur}.
The solver lives in its own compilation unit
(\texttt{GC\_Greedy.hh/cc}), separate from the I/O classes.

% ------------------------------------------------------------
\subsection{The solver loop}
\label{ssec:gc-greedy-loop}
% ------------------------------------------------------------

\begin{lstlisting}[language=C++,
  caption={DSatur main loop (\texttt{GC\_Greedy.cc}).},
  label={lst:gc-greedy-loop}]
void GreedyGCSolver(const Graph& in, Coloring& out) {
  unsigned i, k, c;
  int n;
  // colors[v] = set of colours still available for vertex v
  vector<vector<unsigned>> colors(
      in.Nodes(), vector<unsigned>(in.Colors()));
  for (i = 0; i < in.Nodes(); i++)
    iota(colors[i].begin(), colors[i].end(), 0);

  out.Reset();
  for (i = 0; i < in.Nodes(); i++) {
    n = SelectNode(colors, in, out);  // DSatur selection
    if (colors[n].size() == 0)
      return;  // no feasible colour: search fails
    k = Random(0, colors[n].size() - 1);
    c = colors[n][k];                 // random among available
    Update(n, c, colors, in);         // propagate
    out[n] = c;
  }
}
\end{lstlisting}

\paragraph{Key data structure.}
\texttt{colors[v]} is the list of colours still available for
vertex~$v$.  Initially every vertex has all $k$ colours.
After colouring a vertex, \texttt{Update} removes that colour
from every neighbour's list.
The \emph{saturation degree} of $v$ is implicitly
$k - |\texttt{colors[v]}|$: fewer available colours $=$ higher
saturation.

% ------------------------------------------------------------
\subsection{Node selection (DSatur rule)}
\label{ssec:gc-select-node}
% ------------------------------------------------------------

\begin{lstlisting}[language=C++,
  caption={\texttt{SelectNode}: pick the most saturated
           uncoloured vertex.},
  label={lst:gc-select}]
unsigned SelectNode(const vector<vector<unsigned>>& v,
                    const Graph& in, const Coloring& out) {
  unsigned i;
  int best = -1;
  for (i = 0; i < v.size(); i++)
    if (out[i] == -1 &&               // uncoloured
        (best == -1
         || v[i].size() < v[best].size()    // fewer available
         || (v[i].size() == v[best].size()
             && in.Degree(i) > in.Degree(best))))
      best = i;                        // tie-break: degree
  return best;
}
\end{lstlisting}

\noindent
Smaller \texttt{colors[v].size()} means higher saturation.
Ties broken by highest degree---exactly the DSatur rule.

% ------------------------------------------------------------
\subsection{Propagation}
\label{ssec:gc-update}
% ------------------------------------------------------------

\begin{lstlisting}[language=C++,
  caption={\texttt{Update}: remove colour \texttt{c} from
           all neighbours of~\texttt{n}.},
  label={lst:gc-update}]
void Update(unsigned n, unsigned c,
            vector<vector<unsigned>>& v, const Graph& in) {
  for (unsigned i = 0; i < in.Degree(n); i++) {
    unsigned k = in.ArcListMember(n, i);
    for (unsigned j = 0; j < v[k].size(); j++)
      if (v[k][j] == c) {
        v[k].erase(v[k].begin() + j);
        break;
      }
  }
}
\end{lstlisting}

% ------------------------------------------------------------
\subsection{Driver and Makefile}
\label{ssec:gc-driver}
% ------------------------------------------------------------

\begin{lstlisting}[language=C++,
  caption={Minimal driver (\texttt{GC\_Driver.cc}).},
  label={lst:gc-driver}]
#include "GC_Data.hh"
#include "GC_Greedy.hh"

int main(int argc, char* argv[]) {
  if (argc != 2) {
    cerr << "Usage: " << argv[0] << " <input_file>" << endl;
    return 1;
  }
  Graph in(argv[1]);
  Coloring out(in);
  GreedyGCSolver(in, out);
  cout << out << endl;
  return 0;
}
\end{lstlisting}

\noindent
The driver is trivially short: create Input, create Output, run
Solver, print.  To switch algorithms, only the solver call
changes.

\begin{lstlisting}[language=make,
  caption={Makefile for the graph colouring solver.},
  label={lst:gc-makefile}]
OPTIONS = -Wall -Wfatal-errors -O3

GC_Test.exe: GC_Driver.o GC_Greedy.o GC_Data.o Random.o
    g++ -o GC_Test.exe $^

GC_Data.o: GC_Data.cc GC_Data.hh
    g++ $(OPTIONS) -c GC_Data.cc

GC_Greedy.o: GC_Greedy.cc GC_Greedy.hh GC_Data.hh
    g++ $(OPTIONS) -c GC_Greedy.cc

GC_Driver.o: GC_Driver.cc GC_Data.hh
    g++ $(OPTIONS) -c GC_Driver.cc

clean:
    rm -f *.o GC_Test.exe
\end{lstlisting}


% ============================================================
\section{Example 2: Bus Driver Scheduling (I/O classes)}
\label{sec:bds-data}
% ============================================================

As a second example of the Input/Output pattern, we look at the
\textbf{Bus Driver Scheduling} (BDS) problem: given a set of
\emph{tasks} (bus trips) and a set of \emph{shifts} (driver
work-blocks, each covering certain tasks), select a minimum-cost
set of shifts that covers every task.%
\aside{BDS is a set-covering problem:
shifts = sets, tasks = elements.}

\begin{lstlisting}[language=C++,
  caption={\texttt{BDS\_Data.hh} --- Input class (excerpt).},
  label={lst:bds-input}]
class BDS_Input {
  friend ostream& operator<<(ostream& os, const BDS_Input& bs);
public:
  BDS_Input(string file_name);
  unsigned Tasks()  const { return tasks; }
  unsigned Shifts() const { return shifts; }
  bool Covers(unsigned s, unsigned t) const
    { return covers[s][t]; }
  unsigned ShiftCost(unsigned s) const
    { return shift_cost[s]; }
protected:
  unsigned tasks, shifts;
  vector<unsigned> shift_cost;
  vector<vector<bool>> covers;
  vector<vector<unsigned>> shift_members;
  vector<vector<unsigned>> task_coverage;
};
\end{lstlisting}

\begin{lstlisting}[language=C++,
  caption={\texttt{BDS\_Data.hh} --- Output class (excerpt).},
  label={lst:bds-output}]
class BDS_Output {
  friend ostream& operator<<(ostream& os, const BDS_Output& out);
public:
  BDS_Output(const BDS_Input& i);
  bool operator[](unsigned s) const
    { return selected_shifts[s]; }
  unsigned Uncovered() const { return uncovered; }
  unsigned TotalCost() const { return total_cost; }

  void InsertShift(unsigned s);
  void RemoveShift(unsigned s);
  void RemoveAll();
private:
  const BDS_Input& in;
  vector<bool> selected_shifts;
  vector<unsigned> coverage;
  unsigned uncovered, total_cost;
};
\end{lstlisting}

\paragraph{Design notes.}
\begin{itemize}[nosep]
  \item The Output class maintains \emph{incremental}
    bookkeeping: \texttt{InsertShift} and \texttt{RemoveShift}
    update \texttt{coverage}, \texttt{uncovered}, and
    \texttt{total\_cost} without a full recount.
  \item This incremental pattern will be essential for
    \emph{local search} (later chapters), where we repeatedly
    insert/remove shifts and need fast cost evaluation.
  \item The same Input/Output split applies: greedy,
    backtracking, or local-search solvers all use the same
    \texttt{BDS\_Input} and \texttt{BDS\_Output}.
\end{itemize}
